% Fixed-candidate proofs, included inside Appendix D after its section heading.
\subsection{Fixed endpoints and arbitrary candidates}

\begin{lemma}[Total radial endpoint and maximum-to-own-demand]
\label{lem:appendix-fixed-total-endpoint}
The assertions of \zcref{lem:fixed-total-radial-forcing} hold for arbitrary
V types.
\end{lemma}
\begin{proof}
No closed V triangle contains $O$: its distinguished vertex is interior
and at distance one from $O$, so moving that vertex slightly away from $O$
inside the open triangle would contradict diameter one. Each own trace has
positive length. Hence $0<\gamma_i<1$ and $0<d_i<1$.

If a local open role contained $Z_i(\gamma_i)$, openness would give a point
$Z_i(c)$ with $c>\gamma_i$ in that role. A role with no positive neighboring
trace would acquire one; otherwise its maximal endpoint would increase.
Both contradict the definition of $\gamma_i$. The nonlocal roles are
excluded by
\[
 \|d_iV_i-V_{i\pm2}\|^2=1+d_i+d_i^2>1,
 \qquad \|d_iV_i-V_{i+3}\|=1+d_i>1.
\]
Thus every total endpoint is missed by all open V roles. For
$\gamma_i\le\Gamma<1$, the same proof with $d=1-\Gamma$ excludes
$(1-\Gamma)V_i$: openness would give a local endpoint strictly above
$\Gamma\ge\gamma_i$. For $\Gamma=1$, the point is $O$, which was already
excluded even from the closed V triangles.

Now let an arbitrary open unit candidate contain $O$ and $\widehat P_i$.
If its radial exit is $d_i'$, open containment gives $d_i<d_i'$ and hence
$\gamma_i>1-d_i'$. If both neighboring endpoints are at most $1-d_i'$,
the own term must realize this strict demand. For a nonsupercritical
neighbor with both actual boundary reaches above $d_i'$, the bound
$C_\pm(A_j,B_j)\le1-\min\{A_j,B_j\}<1-d_i'$ excludes its contribution.
No skeleton-cover hypothesis for this new candidate is used.
\end{proof}


\begin{lemma}[Fixed five-point selected-gap enclosure]
\label{lem:appendix-fixed-transverse}
The set $K_{BC}$ in \zcref{prop:new-nplus-one-all-vd0} has
$\Lambda(K_{BC})>1$ for arbitrary V types on the nonsupercritical path,
and is forced into $U_C$ under the stated skeleton-cover hypotheses.
\end{lemma}
\begin{proof}
The boundary-deficit identities give increasing $A_1,\ldots,A_5$ and
decreasing $B_1,\ldots,B_5$. Thus $0<x\le y<1$ and $x/2\le z\le y$.
Put $\rho_i=1-c_{\max}(A_i,B_i)$. Apply
\zcref{lem:new-common-pair-domination} at each neighboring triangle's own pair.
Since $\rho_i\le\min(A_i,B_i)$, path monotonicity bounds the total inward
reach on $r_i$ by
\[
 1-\min\{\rho_i,A_{i-1},B_{i+1}\}\qquad(i=2,3,4).
\]
The bounded-frontier statement of \zcref{lem:fixed-total-radial-forcing}
therefore forces the corresponding points into $U_C$.
The bounds $(A_2,B_2)\ge(1-y,z)$ imply
$r\le\rho_2$, $r\le A_1$, and $r\le B_3$. Hence $rV_2$ is safe.
Similarly $(A_4,B_4)\ge(1-z,x/2)$ gives $t\le\rho_4$, and
$t\le x/2\le B_5$; clipping by $A_3$ makes $\widehat tV_4$ safe.
The gap endpoints follow from \zcref{lem:gap-exhaustion}, and the midpoint
belongs to $U_C$ by hypothesis. All witnesses are fixed before enclosure is tested.

If $A_3\ge t$, apply \zcref{lem:bc-five-point}. Otherwise $A_3<t$ and
$A_3\ge A_1=1-y$, so
\[
 \|X_0(y)-A_3V_4\|^2\ge
 \|X_0(y)-(1-y)V_4\|^2=1+(1-y)(2-y)>1.
\]
The diameter obstruction gives $\Lambda(K_{BC})>1$ in this case also.
No signed center coordinates or CE1 return are used in this enclosure argument.
\end{proof}

\begin{lemma}[Fixed four-point supported-rescuer proof]
\label{lem:appendix-fixed-four-point}
The geometric statement of \zcref{thm:fixed-four-point-rescuer} holds.
\end{lemma}
\begin{proof}
If $a=0$, the set contains $O,V_0$, at distance one. If $s\ge1$, then
\[
 \|Y(a)-\varepsilon V_1\|^2=1-s+s^2\ge1.
\]
These are diameter obstructions. Otherwise put $v=a/s$, so $0<s,v<1$.
The ratio gives $1-\beta\ge v$, while $a=sv\le v$. Hence $Y(v)$ lies
on $[Y(a),Y(1-\beta)]$. By convexity it suffices to enclose the quadrilateral,
in cyclic order,
\[
 O,\qquad Q=Y(v),\qquad A=Y(sv),\qquad P=s(1-v)V_1.
\]
With $h=\sqrt3/2$, its outward edge normals may be chosen as
\[
 (-hv,v/2-1),\quad v(1-s)(h,-1/2),\quad
 (hs,1-s/2),\quad s(1-v)(-h,1/2).
\]
The maximizing triples in each direction and its two $120$-degree rotations
are $(O,A,P)$, $(A,P,O)$, $(A,O,Q)$, and $(O,Q,A)$, respectively.
The resulting caliper sides are
\[
 \frac1{\sqrt{1-v+v^2}},\qquad 1+s(1-v),\qquad
 \frac1{\sqrt{1-s+s^2}},\qquad 1+v(1-s).
\]
All exceed one for $0<s,v<1$. The finite-caliper theorem
\zcref{thm:cert-caliper} completes the nondegenerate case.
\end{proof}
